\section{Non-Localization}


本節では、なぜ意識を特定の場所へ局在化できないのかを整理する。

ここでいう非局在性とは、意識に神経基盤がないという意味ではない。

意識状態は、明らかに身体と脳の活動に依存している。

しかし、その依存関係は、単一部位に意識が「存在する」という形では成立しない。

本稿の立場では、意識は場所ではなく、速度差、結合粘性、内部遅延、意味勾配、ログ参照が一定の関係に入ったときに開かれる時間的レジームである。

したがって、意識を局在化しようとする試みは、意識そのものではなく、意識が通過した後の残余、あるいは意識に関与する局所的基盤を探している可能性が高い。

\subsection{The Mistaken Search for a Seat}

意識研究には、しばしば「意識の座」を探そうとする傾向がある。

どの脳部位が意識を生むのか。

どの回路が主観経験を作るのか。

どの信号が意識そのものに対応するのか。

これらの問いは、一見すると自然である。

しかし、その背後には生成モデルがある。

すなわち、意識はどこかで作られ、その作られたものがどこかに存在する、という前提である。

本稿では、この前提を採らない。

意識は生成物ではなく、開口状態である。

それは、処理とログのあいだに参照可能性が開かれている状態であり、その成立には複数の時間スケール、可塑性ゲート、意味勾配、非閉包ループが関わる。

このため、意識を一点に置こうとする問いは、最初から対象を誤っている。

\subsection{Local Substrates and Non-Local Regimes}

ただし、意識が局在しないということは、意識に局所的基盤がないということではない。

高速層の処理には、それに対応する神経活動がある。

低速層のログ沈殿にも、それに対応する神経的・身体的基盤がある。

報告経路、記憶経路、感覚処理、運動準備、言語化にも、それぞれ局所的な基盤がある。

これらは重要である。

しかし、それらのどれか一つが意識そのものなのではない。

意識は、局所基盤が時間的に結合し、ログ参照可能な関係に入ったときに開かれるレジームである。

したがって、意識研究において必要なのは、局所基盤の否定ではない。

必要なのは、局所基盤と非局在的レジームを混同しないことである。

脳部位は関与する。

しかし、意識状態は部位ではなく、部位間・時間層間・ログ間の関係として成立する。

\subsection{Why Localization Converts Process into Residue}

局在化には、対象を固定する操作が含まれる。

ある現象を特定の場所へ割り当てるためには、その現象を安定した対象として扱う必要がある。

しかし、意識はライブな開口状態であり、完全に固定された対象ではない。

第13節で見たように、測定はクオリアや意識状態をそのまま保持するのではなく、低速層へ沈殿させ、残余へ変換する。

局在化も同じである。

局在化によって得られるのは、ライブな意識そのものではなく、意識状態に関与した残余、痕跡、相関、報告可能なログである。

したがって、ある部位に活動が見られることは重要である。

しかし、それは意識がその部位に存在することを意味しない。

それは、その部位がログ参照ループの一部、またはその残余の一部として関与していることを示す。

局在化は無意味ではない。

ただし、それが捉えるのは意識そのものではなく、意識が開かれたときに関与した局所条件である。

\subsection{Consciousness as Temporal Relation}

本稿の枠組みでは、意識は空間的位置ではなく、時間的関係である。

高速層は未来を先取りする。

低速層は過去を保持する。

可塑性ゲートは、そのあいだで何が沈殿し、何が流れ去るかを調整する。

意味勾配は、どのログが参照されやすいかを方向づける。

この関係が、一定の範囲に入ったとき、意識状態が開かれる。

したがって、意識を問う適切な形式は「どこにあるか」ではない。

「どの時間的条件で、処理がログ参照可能になるか」である。

この条件は、consciousness window として表された。

\begin{equation}
0 < |\Delta v| < \Delta v_{\max}
\end{equation}

\begin{equation}
\mu_{\min} < \mu < \mu_{\max}
\end{equation}

\begin{equation}
\tau_{\min} < \tau < \tau_{\max}
\end{equation}

\begin{equation}
\mathcal{T}_{\min}
<
\mathcal{T}_{\Phi}
<
\mathcal{T}_{\max}
\end{equation}

この窓は空間上の一点ではない。

それは、複数の処理層とログ構造が、時間的に結合している動的領域である。

\subsection{Neural Correlates Are Not Consciousness Itself}

意識の神経相関を探すことは重要である。

しかし、神経相関は意識そのものではない。

神経相関は、意識状態が開かれる条件、または意識状態が開かれた後に残る痕跡を示す。

たとえば、ある神経活動が特定の意識報告と対応している場合、その対応は重要な情報である。

しかし、それだけで、その神経活動が意識そのものであるとは言えない。

なぜなら、同じ活動が、報告条件、注意状態、意味勾配、遅延、可塑性ゲートの状態によって異なる意味を持ちうるからである。

本稿では、神経相関を否定しない。

むしろ、神経相関をログ参照ループの一部として読む。

神経活動は、意識を生成する単独の座ではなく、意識が開かれるための時間的非閉包構造の一要素である。

\subsection{Distributed Does Not Mean Diffuse}

意識が非局在的であると言うと、意識が曖昧に脳全体へ広がっているかのように誤解されることがある。

しかし、本稿の立場はそのような拡散モデルではない。

非局在とは、場所がないという意味ではなく、単一の場所に還元できないという意味である。

意識状態は、複数の局所過程が、特定の時間的関係に入ることで成立する。

したがって、意識はどこにでもあるのではない。

むしろ、かなり厳しい条件のもとでしか開かれない。

速度差が必要である。

粘性が必要である。

遅延が必要である。

意味勾配が必要である。

可塑性ゲートが必要である。

そして、それらが完全閉包にも完全分離にも至らない必要がある。

この意味で、意識は拡散しているのではなく、条件づけられている。

\subsection{From Localization to Conditions}

局在化の問いを完全に捨てる必要はない。

ただし、それは問いの最終形ではない。

「どこが活動しているか」という問いは、「どの局所基盤がログ参照ループに参加しているか」という問いへ変換されるべきである。

さらに、その問いは次のように拡張される。

どの時間スケールが関与しているのか。

どの速度差が維持されているのか。

どの程度の遅延があるのか。

どのログが参照可能になっているのか。

どの意味勾配が参照を方向づけているのか。

どの可塑性ゲートが開閉しているのか。

このように問いを変えることで、意識は「場所探し」の問題ではなく、条件配置の問題になる。

本稿が提案するのは、この問いの移動である。

\subsection{Summary}

本節では、意識の非局在性を整理した。

\begin{itemize}
\item 意識には神経基盤がある。
\item しかし、意識状態そのものは単一部位に局在しない。
\item 局在化が捉えるのは、ライブな意識ではなく、局所基盤や沈殿残余である。
\item 意識は、速度差、結合粘性、内部遅延、意味勾配、可塑性ゲートが作る時間的関係である。
\item 神経相関は重要だが、それは意識そのものではなく、ログ参照ループの一部として解釈されるべきである。
\item 非局在とは、曖昧に広がることではなく、単一部位に還元できないという意味である。
\item 適切な問いは「意識はどこにあるか」ではなく、「どの条件で意識状態が開かれるか」である。
\end{itemize}

したがって、本稿における非局在性は、神秘化ではない。

それは、意識を空間的対象ではなく、時間的非閉包レジームとして扱うことの帰結である。

次節では、このレジームからどのような定性的予測が導かれるのかを整理する。
